%----------------------------------------------------------------------------- %%
\section{Introduction}
\label{bab:1}
%----------------------------------------------------------------------------- %%

Kubernetes schedules Pods by filtering feasible Nodes and scoring the remaining alternatives. This process is effective for initial placement, but running Pods normally remain bound to their Nodes until they terminate or are recreated. As workloads scale and change, clusters accumulate resource fragmentation. Aggregate capacity may be sufficient while no individual Node can accommodate a new Pod. Directed Pod eviction through the Kubernetes Descheduler provides a post-placement correction mechanism by evicting selected Pods and allowing the kube-scheduler to pack them more densely~\cite{k8s_descheduler_docs,larsson_klein_elmroth_2023}.

However, compute-focused consolidation can negatively affect application communication. Moving dependent microservices to distant or degraded Nodes may improve compute packing while significantly increasing latency. Modern microservice architectures are highly sensitive to inter-service communication delays, where a single poor placement decision can cascade into user-facing performance degradation. Communication-aware placement studies show that network topology, observed latency, and service dependencies materially affect placement quality~\cite{marchese2022network,marchese_tomarchio_2022,chen2024trade}. 

Therefore, a practical node reclamation architecture requires a network-aware guardrail. This research develops a Kubernetes Descheduler architecture with a network-aware filtering layer. This layer intercepts eviction decisions made by standard resource-balancing policies, blocking evictions whose feasible destinations would degrade communication locality for dependent services. Our implementation evaluates this trade-off by measuring how effectively a network-aware pre-eviction filter protects communication latency without severely compromising the Descheduler's ability to reduce cluster fragmentation. This study contributes a novel cost filtering mechanism for the Kubernetes Descheduler, demonstrating how strict network-locality preservation interacts with overall cluster resource fragmentation reduction on standard AWS EC2 infrastructure.
